% Style follows ../class2/resume2PT26.tex.
% Build in this directory (run twice for page references):
% pdflatex -interaction=nonstopmode -halt-on-error resume7PT26.tex
% Topics follow the syllabus in ../class1/class1PT26.tex.
% Content adapted from ../lectureNote.tex and ../lectureNotes.tex.
\documentclass[a4paper,10pt,pdflatex,ja=standard,jafont=ipaex-type1,
  magstyle=real,scale=1]{bxjsarticle}
\usepackage{lmodern}
\usepackage{amsmath,amssymb,mathtools}
\usepackage{geometry}
\geometry{reset,left=22mm,right=22mm,top=20mm,bottom=17mm,
  headheight=13pt,headsep=6mm,footskip=10mm}
\usepackage{xcolor}
\usepackage{enumitem}
\usepackage{array,booktabs,tabularx}
\usepackage{tikz}
\usepackage{fancyhdr}
\usepackage{hyperref}
\definecolor{ink}{HTML}{193A5C}
\definecolor{accent}{HTML}{237D82}
\definecolor{soft}{HTML}{EAF3F3}
\definecolor{muted}{HTML}{52616B}
\newcommand{\coursename}{確率論と統計学 I}
\hypersetup{unicode=true,colorlinks=true,urlcolor=accent,linkcolor=ink,
  pdfauthor={中島秀太},pdftitle={第7回：分布と特性関数}}
\pagestyle{fancy}
\fancyhf{}
\fancyhead[L]{\small\color{ink}\coursename{}・授業要約と演習}
\fancyhead[R]{\small\color{ink}第7回}
\fancyfoot[C]{\small\color{ink}\thepage/\pageref*{lastpage}}
\renewcommand{\headrulewidth}{0.3pt}
\setlength{\parindent}{0pt}
\setlength{\parskip}{2pt}
\setlength{\abovedisplayskip}{4pt}
\setlength{\belowdisplayskip}{4pt}
\setlength{\abovedisplayshortskip}{3pt}
\setlength{\belowdisplayshortskip}{4pt}
\setlist[itemize]{leftmargin=1.8em,itemsep=1pt,topsep=2pt,parsep=0pt}
\setlist[enumerate]{leftmargin=2em,itemsep=3pt,topsep=3pt,parsep=0pt}
\newcommand{\topic}[1]{\par\vspace{5pt}
  {\large\sffamily\mdseries\color{accent}#1}\par\vspace{2pt}}
\newcommand{\handouttitle}[1]{%
  \begin{center}
    {\Large\color{ink}#1}\par\vspace{3mm}
    {\color{accent}\rule{0.88\linewidth}{0.7pt}}
  \end{center}\vspace{-2mm}}
\newcommand{\keybox}[1]{\par\vspace{3pt}\begingroup
  \setlength{\fboxsep}{5pt}
  \colorbox{soft}{\parbox{\dimexpr\linewidth-2\fboxsep\relax}{#1}}%
  \endgroup\par\vspace{2pt}}
\newcommand{\examplelabel}[1]{{\sffamily\mdseries\color{ink}#1}}
\newcommand{\columnheading}[1]{\par\vspace{5pt}
  {\fontsize{11pt}{15pt}\selectfont\sffamily\mdseries\color{accent}#1}\par\vspace{2pt}}
\newcommand{\columnstart}[1]{\clearpage
  \setlength{\parskip}{2.5pt}
  \fancyhead[L]{\small\color{ink}\coursename{}・発展的内容}
  {\small\sffamily\mdseries\color{accent}発展的内容}\par\vspace{4pt}
  {\Large\color{ink}#1}\par\vspace{2pt}
  {\color{accent}\rule{\linewidth}{0.7pt}}\par\vspace{4pt}}
\newcommand{\F}{\mathcal F}
\newcommand{\R}{\mathbb R}
\newcommand{\N}{\mathbb N}
\newcommand{\E}{\mathbb E}
\newcommand{\Var}{\operatorname{Var}}
\newcommand{\Cov}{\operatorname{Cov}}
\newcommand{\ind}{\mathrm{1}}
\newcommand{\iid}{\text{独立同分布}}
\begin{document}
\fontsize{10pt}{13.5pt}\selectfont
\begingroup
\fancyhead[L]{\small\color{ink}\coursename{}・前回の演習解答}
\handouttitle{第6回\quad 演習問題の解答}

\topic{演習1：ベルヌーイ和の標準化}
独立な $X_j\sim\mathrm{Bern}(p)$（$0<p<1$）、$S_n=\sum_{j=1}^nX_j$ とする。
$X_j^2=X_j$ より $\E[X_j]=p$、$\Var(X_j)=p(1-p)$。線形性と独立性から
$$
  \E[S_n]=np,\qquad \Var(S_n)=np(1-p),\qquad
  \E\left[\frac{S_n-np}{\sqrt{np(1-p)}}\right]=0,\quad
  \Var\left(\frac{S_n-np}{\sqrt{np(1-p)}}\right)=1.
$$

\topic{演習2：正規近似と連続性補正}
$S_{100}\sim\mathrm{Bin}(100,1/2)$ の平均は $50$、標準偏差は $5$。
標準正規分布関数を $\Phi$ とすると、$P(45\le S_{100}\le55)$ は
$$
  \begin{aligned}
    \text{補正なし：}&\quad\Phi(1)-\Phi(-1)\simeq0.682689,\\
    \text{補正あり：}&\quad\Phi(1.1)-\Phi(-1.1)\simeq0.728668,\\
    \text{厳密値：}&\quad2^{-100}\sum_{k=45}^{55}\binom{100}{k}
      \simeq0.728747.
  \end{aligned}
$$
補正ありでは $[44.5,55.5]$ に対応する正規確率を使い、厳密値に近づいている。

\topic{演習3：一様な局所近似から区間確率へ}
ここでは $p=1/2$ とし、$Z_n=(2S_n-n)/\sqrt n$、$x_{n,k}=(2k-n)/\sqrt n$、
$h_n=2/\sqrt n$、$\phi(x)=e^{-x^2/2}/\sqrt{2\pi}$ とおく。
有限な $a<b$ に対して、局所近似を
$P(S_n=k)=\phi(x_{n,k})h_n(1+r_{n,k})$ と書けば
$\eta_n:=\max_{a\le x_{n,k}\le b}|r_{n,k}|\to0$ であり、
$$
  \left|P(a\le Z_n\le b)-\sum_{a\le x_{n,k}\le b}\phi(x_{n,k})h_n\right|
   \le\eta_n\sum_{a\le x_{n,k}\le b}\phi(x_{n,k})h_n\longrightarrow0.
$$
右の和は Riemann 和として $\int_a^b\phi(x)\,dx$ に収束する。
一様性は、項数が増える和全体の誤差を $\eta_n$ で抑える箇所で使っている。

\topic{演習4：尾の評価から分布関数の収束へ}
$\E[Z_n^2]=1$ なので、Markov の不等式より $P(|Z_n|>R)\le R^{-2}$。
$R>|t|$ とすれば
$$
  |P(Z_n\le t)-\Phi(t)|
    \le\left|P(-R\le Z_n\le t)-\int_{-R}^t\phi(x)\,dx\right|
       +R^{-2}+\Phi(-R).
$$
まず $n\to\infty$ とすると第1項は消え、次に $R\to\infty$ とすれば残りも消える。
したがって、すべての $t\in\R$ で $P(Z_n\le t)\to\Phi(t)$ となる。
\clearpage
\endgroup

\handouttitle{第7回\quad 分布と特性関数}
\keybox{分布を振動関数の期待値で表すと、独立な和は積になる。
特性関数は常に存在し、分布の一意性と極限の判定に使える。}

\topic{特性関数の定義と基本性質}
実数値確率変数 $X$ の特性関数を、$i^2=-1$ として
$$\varphi_X(t):=\E[e^{itX}]=\int_\R e^{itx}\,\mu_X(dx)
 \qquad(t\in\R)$$
と定める。複素数値の期待値は実部と虚部に分けて定義する。
$|e^{itX}|=1$ より、$X$ が可積分でなくても存在し、
$$\varphi_X(0)=1,\qquad |\varphi_X(t)|\le1,\qquad
 \varphi_X(-t)=\overline{\varphi_X(t)}.$$
$t_n\to t$ なら、優収束定理によって $\varphi_X(t_n)\to\varphi_X(t)$。
また $a,b\in\R$ に対し $\varphi_{aX+b}(t)=e^{itb}\varphi_X(at)$ である。

\topic{代表的な分布の特性関数}
\begin{center}
\renewcommand{\arraystretch}{1.6}
\begin{tabular}{@{}ll@{}}
\toprule
分布 & 特性関数 $\varphi_X(t)$ \\
\midrule
$\mathrm{Bern}(p)$ & $1-p+pe^{it}$ \\
$\mathrm{Bin}(n,p)$ & $(1-p+pe^{it})^n$ \\
$\mathrm{Poi}(\lambda)$ & $\exp(\lambda(e^{it}-1))$ \\
$\mathrm{Exp}(\lambda)$ & $\lambda/(\lambda-it)$ \\
$N(m,\sigma^2)$ & $\exp(imt-\sigma^2t^2/2)$ \\
\bottomrule
\end{tabular}
\end{center}
たとえばポアソン分布では、指数関数の級数から
$$\varphi_X(t)=e^{-\lambda}\sum_{k=0}^\infty
 \frac{(\lambda e^{it})^k}{k!}=\exp(\lambda(e^{it}-1)).$$

\topic{独立な和を積に直す}
独立な $X,Y$ に対し、$e^{itX},e^{itY}$ も独立で有界なので
$$\varphi_{X+Y}(t)=\E[e^{itX}e^{itY}]
 =\varphi_X(t)\varphi_Y(t).$$
独立な $X\sim\mathrm{Poi}(\lambda)$、$Y\sim\mathrm{Poi}(\mu)$ なら
$$\varphi_{X+Y}(t)=\exp((\lambda+\mu)(e^{it}-1))$$
となる。次ページの一意性定理により、$X+Y\sim\mathrm{Poi}(\lambda+\mu)$ と分かる。
和の分布の計算を、関数の積と既知の形の比較へ置き換えられる。

\newpage
\topic{Fourier 変換と Gauss 核による正則化}
可積分な関数 $f$ の Fourier 変換を
$$\widehat f(t):=\int_\R f(x)e^{-itx}\,dx$$
と定める。$f\in C_c(\R)$（コンパクトな台をもつ連続関数）でも、
$\widehat f$ が可積分とは限らない。そのため逆変換には条件が必要となる。
$$q_\varepsilon(x)=\frac1{\sqrt{2\pi}\varepsilon}
 e^{-x^2/(2\varepsilon^2)},\qquad
 \widehat q_\varepsilon(t)=e^{-\varepsilon^2t^2/2}\quad(\varepsilon>0)$$
を使って平滑化すると、Fubini の定理により
$$
 (f*q_\varepsilon)(x)
 =\frac1{2\pi}\int_\R\widehat f(t)e^{-\varepsilon^2t^2/2}e^{itx}\,dt,
 \qquad (f*q_\varepsilon)(x):=\int_\R f(x-y)q_\varepsilon(y)\,dy.
$$
$\varepsilon\downarrow0$ で $f*q_\varepsilon\to f$ は一様収束する。
$\widehat f\in L^1$ も仮定すれば、優収束定理により Gauss 因子を外せる。

\topic{特性関数は分布を一意に決める}
上の式を期待値に入れると
$$\E[(f*q_\varepsilon)(X)]
 =\frac1{2\pi}\int_\R\widehat f(t)e^{-\varepsilon^2t^2/2}\varphi_X(t)\,dt.$$
$\varphi_X=\varphi_Y$ なら平滑化した期待値が一致し、$\varepsilon\downarrow0$ で
$\E[f(X)]=\E[f(Y)]$（すべての $f\in C_c(\R)$）となる。
この積分の一致は $\mu_X=\mu_Y$ を意味する。
開区間の指示関数を下から連続関数で近似し、測度の一意性を使えばよい。

\topic{特性関数による収束の判定}
分布収束 $X_n\Rightarrow X$ とは、任意の有界連続関数 $f$ に対して
$\E[f(X_n)]\to\E[f(X)]$ となることである。
L\'evy の連続性定理は
$$X_n\Rightarrow X\quad\Longleftrightarrow\quad
 \varphi_{X_n}(t)\to\varphi_X(t)\quad(\forall t\in\R)$$
を与える。さらに点ごとの極限 $\varphi$ が原点で連続なら、それ自身が
ある確率分布の特性関数となる。分布収束とこの条件の役割は第9回で詳しく扱う。

\topic{モーメントが有限なら原点で微分できる}
$\E|X|<\infty$ なら $\varphi_X'(0)=i\E[X]$。
$\E[X^2]<\infty$ なら $\varphi_X''(0)=-\E[X^2]$ であり、
$$\varphi_X(t)=1+it\E[X]-\frac{t^2}2\E[X^2]+o(t^2).$$
この展開と独立な和の積公式が、一般の中心極限定理の証明の出発点になる。

\newpage
\topic{演習問題}
\begin{enumerate}[label=\arabic*.]
\item $X\sim\mathrm{Geo}(p)$（$P(X=k)=p(1-p)^{k-1}$、$k\ge1$）の特性関数を求めよ。
得られた分母が実数 $t$ について消えないことを確かめよ。
\item $X\sim\mathrm U(-a,a)$（$a>0$）の特性関数を求め、$t=0$ における値と連続性を示せ。
\item 独立な $X\sim N(m_1,v_1)$、$Y\sim N(m_2,v_2)$（$v_1,v_2>0$）について、
特性関数から和の分布を求めよ。
\item $f\in C_c(\R)$ とする。任意の $\delta>0$ に対して
$$\|f*q_\varepsilon-f\|_\infty
 \le\sup_{|u-v|\le\delta}|f(u)-f(v)|
 +2\|f\|_\infty\int_{|y|>\delta}q_\varepsilon(y)\,dy$$
を示し、一様収束を導け。
\end{enumerate}

\topic{レポート問題}
\begin{enumerate}[label=\arabic*.]
\item $\mathrm{Gamma}(m,\beta)$（整数 $m\ge1$、rate $\beta>0$）の特性関数を、
独立な指数分布の和を用いて求めよ。
\item Gauss 核による正則化逆変換公式を証明せよ。
$q_\varepsilon$ の逆表示を畳み込みに代入し、積分順序の交換が可能なことを
絶対値の積分で示すこと。続いて $\widehat f\in L^1$ の場合に古典的な逆変換公式を導け。
\end{enumerate}

\topic{考えるための手がかり}
$\widehat f$ の符号は $e^{-itx}$、特性関数の符号は $e^{itX}$ としている。
Fourier 変換の定義を途中で変えないようにしよう。
正則化では $|\widehat f(t)|\le\|f\|_1$ と Gauss 因子の可積分性が使える。
\vfill
{\small\color{muted}解答メモ：逆変換の存在、Fubini の条件、極限操作の条件をそれぞれ明記する。}

\columnstart{平均がなくても、分布の波形は記述できる}
期待値や分散は確率分布の要約として便利だが、有限とは限らない。
特性関数は、値の大きさを振動に変えてから平均することで、どの分布にも定義できる。

\columnheading{振動を平均する}
Euler の公式から
$$\varphi_X(t)=\E[\cos(tX)]+i\E[\sin(tX)].$$
周波数 $t$ ごとに、$X$ を角度 $tX$ へ写し、単位円上の点を平均していると考えられる。
観測値が大きくなっても単位円の外へは出ないため、可積分性の問題が起こらない。
ただし、振動の打ち消し合いも含めて分布を表すので、単なる平均値の一覧ではない。

\columnheading{積率母関数との違い}
積率母関数 $M_X(t)=\E[e^{tX}]$ では、$e^{tX}$ が片方の尾で大きくなる。
たとえば標準 Cauchy 分布は密度
$$f(x)=\frac1{\pi(1+x^2)}$$
をもち、$\E|X|=\infty$、$M_X(t)=\infty$（$t\ne0$）である。
一方、その特性関数は $\varphi_X(t)=e^{-|t|}$ となる。
原点で微分できないことにも、有限の一次モーメントをもたない性質が現れている。

\columnheading{一意性と、有限個のモーメントの違い}
平均と分散が一致する異なる分布は容易に作れる。
たとえば公平な符号 $R\in\{-1,1\}$ と $G\sim N(0,1)$ はともに平均 $0$、分散 $1$ だが、
$$\varphi_R(t)=\cos t,\qquad \varphi_G(t)=e^{-t^2/2}$$
は異なる。特性関数をすべての周波数で知ることは、
原点での少数の微分係数だけを知ることより多くの情報を含む。
一意性定理は、その情報が分布全体を決定することを保証する。

\columnheading{和の極限で二次の項が残る}
平均 $0$、分散 $1$ の独立同分布列では、原点近くの展開から
$$\varphi_X(u)=1-\frac{u^2}2+o(u^2),\qquad
 \varphi_{n^{-1/2}\sum_{j=1}^nX_j}(t)
 =\left(1-\frac{t^2}{2n}+o(n^{-1})\right)^n\longrightarrow e^{-t^2/2}.$$
L\'evy の定理によって正規分布への収束が従う。
複雑な畳み込みを何度も計算せず、原点近くの形を読むだけで和の極限を捉えられる。
この方法が働くために二次モーメントをどこで使ったかは、第10回で確かめよう。

\label{lastpage}
\end{document}
